% Generated from validated Article Draft Result. Do not hand-edit; revise the structured draft instead.

2024年上期にはLlama 3、OLMo、DeepSeek-V2、DBRX、Mixtral 8x22B、Qwen2、Gemmaなど、公開形態の異なるmodel familyが相次いだ。これらを単純に「openかclosedか」で二分すると、重みだけが公開される場合と、学習コード・データ・中間成果物まで追跡できる場合の差が消える。openは二値属性ではなく、再現可能性を構成する複数の層として捉える必要がある。\autocite{src-0395647db4a8,src-05f98c47f3e8,src-3d2fa83e273d,src-461ae12066d1,src-d9d69089eb6a,src-f9616f30e8f5,src-ff36b1c2ef30}


OLMoが象徴したのは、open weightsより一段深いopen trainingの志向である。モデルそのものに加えて学習過程を研究対象として扱えることは、結果を利用するだけでなく、なぜその結果になったかを検証する余地を増やす。OpenELMもtraining/inference frameworkを含む公開を掲げ、小型モデル側から同じ問いを投げかけた。\autocite{src-3d2fa83e273d,src-4d404eabda6e}


一方、Llama 3は広い配備基盤を持つopen-weight familyとして存在感を高め、DeepSeek-V2やDBRX、Mixtral 8x22BはMoEを含む効率設計を前面に出した。Qwen2やGemmaも異なる規模・用途の選択肢を広げた。ここで競争しているのは単なるparameter countではなく、利用可能な計算資源に対してどの程度のモデルを配置できるかというdeployment economicsでもある。\autocite{src-0395647db4a8,src-05f98c47f3e8,src-461ae12066d1,src-d9d69089eb6a,src-f9616f30e8f5,src-ff36b1c2ef30}


\begin{claimboundary}
「open」という表現は、license、学習データの開示、training recipe、checkpoint、評価再現性を自動的に保証しない。各projectが示した性能・効率の数値も、それぞれの一次資料に帰属する。本稿では公開範囲の違いを比較するが、公開されたweightだけから完全な学習再現性を推定しない。\autocite{src-37f17fc6592f,src-3d2fa83e273d,src-8ccec4bc715c,src-d9d69089eb6a,src-f9616f30e8f5}
\end{claimboundary}

この半年に見えてきたのは、open modelの価値が「無料でdownloadできる」ことだけではなく、研究・fine-tuning・量子化・local serving・独自deploymentへどれだけ接続できるかで決まるという構図である。2024-H1は、open-weight競争がmodel stack競争へ変わり始めた時期として位置付けられる。\autocite{src-3d2fa83e273d,src-4d404eabda6e,src-d9d69089eb6a,src-f9616f30e8f5}
